\documentclass[11pt]{letter}
\usepackage[a4paper,left=2.5cm, right=2.5cm, top=1cm, bottom=1cm]{geometry}
\usepackage{hyperref}
\usepackage[osf]{mathpazo}
\signature{Caleb Scutt}
\address{School of Biosciences\\University of Sheffield\\Sheffield, S10 2TN\\United Kingdom\\cnscutt1@sheffield.ac.uk}
\longindentation=0pt
\begin{document}

\begin{letter}{}
\opening{Dear Editors,}

Since the seminal publications by Foote (1997) and Wills (1994), morphological disparity has become a near ubiquitous metric used in macroevolutionary studies.
Measures of disparity through time and between clades has been used extensively for answering questions about selectivity across mass extinctions, convergence and interspecific competition.
However, these analyses are inevitably limited by the incomplete fossil record, which can bias our inferences from disparity.
Ancestral state estimation (ASE) has been proposed as a way `filling in the gaps' left by an incomplete fossil record, allowing higher resolution of disparity through time analyses for example.
However, the precise impact of methodological choices, such as whether ASE should be conducted before or after trait space ordination and the handling of estimation uncertainty, on the accurate recovery of trait space remains under-explored.

Here in \textit{``Disparity analyses are robust to ancestral state estimation uncertainty''}, we provide guidelines for using ancestral state estimations in disparity analyses, provided through robust benchmarking of methods through simulations.
We show that, for continuous characters, probabilistic approaches (propagating uncertainty) to ancestral state estimation yield more accurate disparity estimates than tip-only analyses, but highlight caution that model misspecification and poor fossil sampling can lead to erroneous inferences when ancestral states are included.
For discrete characters, we find a contextual pattern where the optimal method of ASE is dependent on transition rate and fossil sampling, but, crucially, where the inclusion of ancestral states improves estimates of disparity over their exclusion consistently.
Ultimately, however, we find that fossil sampling density to be the strongest driver of estimation accuracy, which serves as a stark reminder for their inclusion in macroevolutionary workflows.

We have structured this manuscript to serve as a comprehensive and highly practical guide to ASE in disparity analyses.
Given the widespread use of these metrics, we believe our evidence-based recommendations will be of immediate and significant value to the diverse readership of \textit{Palaeontology}.


% Since the seminal publication by Petchey and Gaston in Ecology Letters in 2002 \textit{``Functional diversity (FD), species richness and community composition''} (cumulating $>$ 2000 citations), functional diversity is now routinely used as a statistic to describe complex multidimensional datasets.
% This has been driven by a really motivated community of ecologists proposing many different ways to measure functional diversity depending on different datasets or different questions.
% More than 20 years later however, despite the enthusiasm of the ecology and evolution community there is no consensual way on which metric to use and in which scenarios they perform better.
% This has led to the positive and rich paradigm of a diverse number of metrics routinely used for different tasks.

% However, this richness of ways to measure functional diversity can make it daunting and or confusing for researchers that want to add functional diversity to their research methodology toolkit.
% We argue that this confusions can stem from both (1) an epistemological misunderstanding as well as (2) practical hurdles.
% The latter is nicely alleviated by a dynamic community of researchers proposing open source and easy to use implementations to measure functional diveristy.
% But the former misunderstanding remains.
% Indeed, although it is often rather clear how functional diversity in general can capture changes communities (investigating processes; answering the question ``how are communities changing?'') and relate them to general understanding in ecology (understanding mechanisms; ``why are communities changing?''); surprisingly little work has been published in the last few years on how the metrics used to measure functional diversity can accurately measure complex multidimensional patterns (capturing patterns; ``what is changing?'').

% Here, in \textit{``The what, how and why of trait-based analyses in ecology''} we argue that a better understanding of the ``capturability'' of a pattern should be taken more into account for functional diversity to move forward as a great tool in ecology and evolution.
% Through simulated and empirical data we show that the choice of the functional diversity metric (what) can lead to often very different interpretations of the same processes and mechanisms (how and why).
% For example, in our empirical data looking at endemic birds extinctions in Hawaii we show that different functional diversity metrics can lead to widely different interpretations of the actual biological mechanisms changing Hawaiian bird communities.
% I.e. with the same dataset and a drastic rate of extinction for endemic birds, different metrics can either show a clear decrease in functional diversity (expected) but also a clear increase (unexpected).
% We have written this manuscript with an emphasis on the didactic side of the question and we hope it will be of great interest and usefulness to the many readers of Ecology Letters using functional diversity in their research.

% I look forward to hearing from you soon.

\closing{Yours sincerely,\\On behalf of my co-authors}

\end{letter}
\end{document}

% Editors for Ecology letters:
% David Mouillot, France
% email: david.mouillot@univ-montp2.fr
% Fish, functional diversity, macroecology, model, indices

% Nathan Swenson, USA
% email: swenson@umd.edu
% Functional trait, phylogeny, tropical forest ecology, forest dynamics plot, niche conservatism, trait evolution, transcriptome, functional ecology, macro ecology, community ecology

% Rob Salguero-Gómez, UK
% email: rob.salguero@zoo.ox.ac.uk
% Functional traits, Demography, Phylogenetic analyses, Comparative biology, Life history evolution, Senescence, Integral Projection Models, Matrix Population Models, Phylogenetic comparisons

% Cyrille Violle, France
% email: cyrille.violle@cefe.cnrs.fr
% Functional ecology, functional biogeography, community ecology, macroecology, species coexistence, allometry, domestication


% Suggested reviewers

% Alejandro Ordonez
% Aarhus University
% alejandro.ordonez@bio.au.dk

% Enrico Tordoni
% University of Tartu 
% enrico.tordoni@ut.ee

% Aurele Toussaint 
% Universite Toulouse III-Paul Sabatier 
% aurele.toussaint@univ-tlse3.fr 

% Simona Bonelli
% Università di Torino 
% simona.bonelli@unito.it